\documentclass[11pt]{article}

\usepackage[utf8]{inputenc}
\usepackage[T1]{fontenc}
\usepackage{lmodern}
\usepackage{geometry}
\usepackage{amsmath,amssymb,amsfonts,amsthm,mathtools}
\usepackage{bm}
\usepackage{enumitem}
\usepackage{microtype}
\usepackage{hyperref}
\usepackage{titlesec}
\usepackage{booktabs}
\usepackage{array}
\usepackage{setspace}
\usepackage{mathrsfs}
\usepackage{physics}

\geometry{margin=1in}
\hypersetup{colorlinks=true, linkcolor=black, urlcolor=blue, citecolor=black}
\setlist[enumerate]{topsep=0.4em,itemsep=0.35em}
\setlist[itemize]{topsep=0.3em,itemsep=0.25em}
\titleformat{\section}{\large\bfseries}{\thesection.}{0.5em}{}
\titleformat{\subsection}{\normalsize\bfseries}{\thesubsection.}{0.5em}{}
\setstretch{1.06}

\newcommand{\Aether}{\AE{}ther}
\newcommand{\Aflow}{\AE{}ther-flow}
\newcommand{\TheoryName}{\Aflow{} Interpretation of Relativity}
\newcommand{\TheoryProgram}{\Aether{} / \Aflow{} framework}
\newcommand{\ObserverMetric}{g^{\mathrm{br}}}
\newcommand{\CandidateMetric}{g^{\mathrm{cand}}}
\newcommand{\CompletedMetric}{g^{\sharp}}
\newcommand{\RedefMetric}{\widehat g}
\newcommand{\BaseShift}{\Sigma^{\mathrm{IER}}}
\newcommand{\ResponseShift}{\Sigma^{\mathrm{resp}}}
\newcommand{\CompShift}{\Sigma^{\mathrm{cmp}}}
\newcommand{\ResidualCore}{\mathcal V_{\mu\nu}^{\mathrm{res}}}
\newcommand{\ResidualDec}{\mathcal D_{\mu\nu}^{\mathrm{res}}}
\newcommand{\MixQuad}{\mathcal M_{\mu\nu}^{\mathrm{mix}}}
\newcommand{\RespQuad}{\mathcal Q_{\mu\nu}^{\mathrm{resp}}}
\newcommand{\DeltaTCand}{\Delta T_{\mu\nu}^{\mathrm{cand}}}
\newcommand{\ReducedEin}{\mathcal L_{\ObserverMetric}}
\newcommand{\GaugeBook}{\mathcal G_{\ObserverMetric}}
\newcommand{\CompMix}{\mathcal M_{\mu\nu}^{\mathrm{cmp}}}
\newcommand{\CompQuad}{\mathcal Q_{\mu\nu}^{\mathrm{cmp}}}
\newcommand{\DeltaTComp}{\Delta T_{\mu\nu}^{\mathrm{cmp}}}
\newcommand{\ObsBurden}{\mathfrak B_{\mu\nu}^{\mathrm{obs}}}
\newcommand{\ObserverClass}[1]{\left[#1\right]_{\mathrm{obs}}}

\newtheorem{definition}{Definition}
\newtheorem{proposition}{Proposition}
\newtheorem{remark}{Remark}
\newtheorem{corollary}{Corollary}
\newtheorem{theorem}{Theorem}

\title{\textbf{The \TheoryName{}}\\[0.4em]
\large Substrate-Response / Self-Response Charge-Polarization Candidate\\
\large Exact-Preservation Observer-Equation Same-Package Field-Redefinition Test Note}
\author{}
\date{}

\begin{document}

\maketitle

\begin{abstract}
The active one-metric exact-preservation branch has already passed the first ordered closure-upgrade test. Direct exact absorbability of
\[
\ObsBurden
=
\ResidualDec+\CompMix+\CompQuad-8\pi G_N\,\DeltaTComp
\]
has now been tested on the fixed package and is not yet established. The second ordered test is therefore the only remaining same-package closure route on that branch: can one exhibit an explicit exact invertible field redefinition of the completed observer equation into Einstein--Hilbert plus ordinary matter through one operative metric only, with no second matter law and modulo only the already-extracted gauge-exact and constraint-exact sectors?

The present manuscript performs that second test on the currently recorded fixed package. It proves first that a mere relabeling of the completed metric does not count as a successful field-redefinition upgrade: if \(\RedefMetric\) is only \(\CompletedMetric\) under another name, then the rewritten Einstein--matter quotient statement collapses back to the direct absorbability problem already tested. It then records the decisive package-level fact. The current fixed package contains exact non-independence and same-package re-expansion results, but it does not yet contain an explicit invertible rewritten metric variable together with a proved Einstein--matter reduction theorem of the required type.

Accordingly, the same-package field-redefinition test does not yet close on the current package. Combining this with the already-recorded direct absorbability verdict, neither ordered upgrade route is presently established on the fixed one-metric branch. The strongest proved exact observer-side theorem therefore remains effective equivalence rather than strict observer-side closure, and the next honest primary paper on that same fixed bottleneck is now a genuine necessity, maximality, or no-go theorem. A deeper-origin continuation is primary only if it directly supplies the missing stronger absorbability identity or an explicit same-package field-redefinition theorem.
\end{abstract}

\section{Introduction}

The observer-equation sequence on the active one-metric exact-preservation branch has now reached a sharper decision point than before. The completed observer equation was already reduced to one displayed burden \(\ObsBurden\). That burden was then classified sector by sector, and afterwards the displayed sectors were shown not to survive as additional independent observer-side physics. The first ordered upgrade test was then executed directly on the exact observer quotient, and it did not yet prove \(\ObserverClass{\ObsBurden}=0\).

That leaves only one same-package closure route on the fixed branch: an explicit exact invertible field redefinition that rewrites the completed observer equation itself into Einstein--Hilbert plus ordinary matter through one operative metric only, with no second matter law and modulo only the already-extracted gauge-exact and constraint-exact sectors.

The present note performs that second ordered test and nothing else. It does not add a deeper-origin layer, and it does not claim a no-go theorem by itself. Its narrower task is to decide what the currently recorded fixed package does and does not already prove about the field-redefinition route, and then to state the next honest paper consequence of that verdict.

\paragraph{Framework and claim boundary.}
\TheoryProgram{} denotes the broader \Aether{} / \Aflow{} framework built on the claims that the \Aether{} is the underlying four-dimensional substrate of reality, the \Aflow{} is its intrinsic ordered motion, observed three-dimensional space is the local experiential slice of that deeper substrate, S-time is the experienced order of change arising from matter, light, and the \Aflow{}, and observed expansion is the three-dimensional appearance of deeper four-dimensional ordered motion. Within that framework, \TheoryName{} denotes the exact-closure theory in which the effective gravitational dynamics are adopted to be exactly Einsteinian with universal matter coupling. In the active sequence, \emph{adoption} means use of that established relativistic dynamics without claiming substrate derivation, while \emph{derivation} is reserved for a first-principles recovery from explicit substrate variables.

The present note stays strictly inside that boundary. It does not claim a completed first-principles derivation of GR from substrate microphysics. It does not claim the strict observer-side closure theorem. It does not claim a universal impossibility theorem against all future enlarged packages. Its positive role is narrower: execute the second ordered exact test on the currently fixed package and record the next honest primary-paper consequence of that test.

\section{Fixed Package and Second-Test Criterion}

\begin{definition}[Fixed one-metric exact-preservation package]
\label{def:fixed}
The present note keeps the following observer-level data fixed.
\begin{enumerate}
    \item The candidate and completed operative metrics are
    \begin{equation}
    \CandidateMetric
    =
    \ObserverMetric+\BaseShift+\ResponseShift,
    \qquad
    \CompletedMetric
    =
    \ObserverMetric+\BaseShift+\ResponseShift+\CompShift.
    \label{eq:metrics}
    \end{equation}
    \item The fixed completion solve is
    \begin{equation}
    \ReducedEin(\CompShift)=-\ResidualCore.
    \label{eq:solve}
    \end{equation}
    \item The completed observer equation is
    \begin{equation}
    G_{\mu\nu}[\CompletedMetric]
    =
    8\pi G_N\,T_{\mu\nu}^{\mathrm{matter}}[\CompletedMetric,\psi]
    +\GaugeBook(\CompShift)_{\mu\nu}
    +\ObsBurden,
    \label{eq:completed}
    \end{equation}
    with
    \begin{equation}
    \ObsBurden
    \equiv
    \ResidualDec+\CompMix+\CompQuad-8\pi G_N\,\DeltaTComp.
    \label{eq:burden}
    \end{equation}
    \item The inherited candidate-branch decomposition
    \begin{equation}
    \ResidualDec
    =
    \MixQuad+\RespQuad-8\pi G_N\,\DeltaTCand
    \label{eq:resdec}
    \end{equation}
    is kept fixed.
    \item Gauge bookkeeping is already extracted explicitly in \eqref{eq:completed}, and the constraint-exact elimination sectors have already been removed before \(\ObsBurden\) is defined.
    \item The strongest exact observer-side theorem already proved on this branch is exact effective equivalence: the displayed sectors in \(\ObsBurden\) are exactly non-independent, so the completed branch is exactly observer-side effectively equivalent to Einstein--Hilbert plus ordinary matter through the same one operative metric \(\CompletedMetric\).
    \item The first ordered closure-upgrade test has already been executed on this same fixed package, and direct exact absorbability of \(\ObsBurden\) is not yet established there.
\end{enumerate}
\end{definition}

\begin{definition}[Exact observer quotient]
\label{def:quot}
For the fixed package of Definition \ref{def:fixed}, the \emph{exact observer quotient} is the quotient of observer-side tensors by the already-extracted null classes:
\begin{enumerate}
    \item gauge-exact bookkeeping tensors;
    \item constraint-exact elimination tensors.
\end{enumerate}
The observer-quotient class of a tensor \(X_{\mu\nu}\) is denoted \(\ObserverClass{X_{\mu\nu}}\).
\end{definition}

\begin{definition}[Explicit same-package field-redefinition upgrade]
\label{def:redef}
An \emph{explicit same-package field-redefinition upgrade} on the fixed package means an explicit theorem exhibiting an exact invertible rewriting of the completed observer equation in terms of a rewritten metric variable \(\RedefMetric_{\mu\nu}\) such that:
\begin{enumerate}
    \item \(\RedefMetric_{\mu\nu}\) is determined only by the already fixed branch data;
    \item no second operative metric and no second matter law are introduced;
    \item the rewritten observer equation is exactly Einstein--Hilbert plus ordinary matter through that same rewritten one-metric route, modulo only the already-extracted gauge-exact and constraint-exact sectors.
\end{enumerate}
\end{definition}

\begin{remark}
The second ordered test is stronger than same-package origin language. To pass Definition \ref{def:redef}, one needs an explicit invertible rewritten metric variable and a proved Einstein--matter reduction theorem in that rewritten variable, not merely the observation that the displayed sectors come from the same one-metric re-expansion package.
\end{remark}

\section{What a Genuine Field-Redefinition Upgrade Must Add}

\begin{proposition}[A trivial relabeling of \texorpdfstring{\(\CompletedMetric\)}{g-sharp} does not create a new upgrade route]
\label{prop:trivial}
Suppose one sets
\[
\RedefMetric_{\mu\nu}=\CompletedMetric_{\mu\nu}
\]
without adding any new exact identity beyond the already recorded fixed package. Then a successful rewritten Einstein--matter quotient statement is equivalent to the direct absorbability statement already tested on that same package.
\end{proposition}

\begin{proof}
With \(\RedefMetric=\CompletedMetric\), the rewritten observer equation is just the original completed observer equation \eqref{eq:completed} under another name. Therefore an exact quotient identity of the form
\[
\ObserverClass{G_{\mu\nu}[\RedefMetric]-8\pi G_N\,T_{\mu\nu}^{\mathrm{matter}}[\RedefMetric,\psi]}=0
\]
becomes
\[
\ObserverClass{G_{\mu\nu}[\CompletedMetric]-8\pi G_N\,T_{\mu\nu}^{\mathrm{matter}}[\CompletedMetric,\psi]}=0.
\]
By \eqref{eq:completed}, that is exactly the statement that the observer-quotient class of \(\ObsBurden\) vanishes after the already-extracted gauge-exact bookkeeping sector is quotiented away. In other words, the trivial relabeling route collapses back to the direct absorbability problem rather than furnishing a new independent upgrade route.
\end{proof}

\begin{remark}
Proposition \ref{prop:trivial} explains why the second test cannot be passed merely by renaming \(g^\sharp\). A genuine success on the field-redefinition route must supply new explicit rewriting data beyond the already failed direct absorbability presentation.
\end{remark}

\begin{proposition}[What the current fixed package already gives]
\label{prop:current}
On the current fixed package:
\begin{enumerate}
    \item the displayed sectors in \(\ObsBurden\) are exact same-package re-expansion sectors and are exactly non-independent;
    \item the same package does not yet display an explicit invertible rewritten metric variable \(\RedefMetric_{\mu\nu}\) together with a proved Einstein--matter reduction theorem of the type required by Definition \ref{def:redef}.
\end{enumerate}
\end{proposition}

\begin{proof}
Item (1) is part of Definition \ref{def:fixed}: the prior observer-equation notes already prove exact non-independence and exact same-package origin for \(\ResidualDec\), \(\CompMix\), \(\CompQuad\), and \(\DeltaTComp\). For item (2), inspect the currently recorded fixed package itself. Its exact observer-side results are the completed equation \eqref{eq:completed}, the exact non-independence / effective-equivalence theorem, and the direct absorbability verdict. None of those recorded results exhibits an explicit invertible rewritten metric variable \(\RedefMetric_{\mu\nu}\) or proves that the completed observer equation becomes exactly Einstein--Hilbert plus ordinary matter in that rewritten variable modulo only the already-extracted null sectors. Therefore the present package does not yet contain the required field-redefinition theorem.
\end{proof}

\section{Same-Package Field-Redefinition Test on the Current Package}

\begin{theorem}[Current verdict of the same-package field-redefinition test]
\label{thm:test}
On the currently recorded fixed package, an explicit same-package field-redefinition upgrade is not yet established.
\end{theorem}

\begin{proof}
Definition \ref{def:redef} requires an explicit invertible rewritten metric variable and a proved Einstein--matter reduction theorem in that rewritten variable. Proposition \ref{prop:trivial} shows that a trivial relabeling of the completed metric does not provide a new route at all; it simply collapses back to the direct absorbability problem already tested and not yet established on the same package. Proposition \ref{prop:current}(2) then gives the decisive current-status fact: the recorded fixed package still does not contain any nontrivial explicit invertible rewritten metric variable together with the required exact Einstein--matter reduction theorem. Therefore the same-package field-redefinition upgrade is not yet established on the current package.
\end{proof}

\begin{corollary}[Current strongest exact observer-side theorem after both ordered tests]
\label{cor:strongest}
After the direct absorbability test and the same-package field-redefinition test, the strongest exact observer-side theorem still proved on the current one-metric branch remains the exact observer-side effective-equivalence theorem rather than strict observer-side closure.
\end{corollary}

\begin{proof}
Definition \ref{def:fixed} records that effective equivalence is already proved on the completed branch and that direct exact absorbability is not yet established there. Theorem \ref{thm:test} adds that the remaining same-package field-redefinition route is also not yet established on the current fixed package. Therefore neither closure-upgrade route has yet been proved, and effective equivalence remains the strongest exact observer-side theorem currently available.
\end{proof}

\section{Consequence for Project Priority}

\begin{theorem}[Next honest primary paper on the fixed bottleneck]
\label{thm:next}
Because neither direct exact absorbability nor an explicit same-package field-redefinition upgrade is presently established on the current fixed one-metric package, the next honest primary paper on that same observer-side bottleneck is a genuine necessity, maximality, or no-go theorem.
\end{theorem}

\begin{proof}
Item (7) of Definition \ref{def:fixed} records the outcome of the first ordered test: direct exact absorbability is not yet established on the current package. Theorem \ref{thm:test} gives the outcome of the second ordered test: the explicit same-package field-redefinition route is also not yet established on that same package. Those are exactly the two closure-upgrade routes isolated earlier for the fixed branch. Since the current package has now been carried through both tests without establishing either route, another manuscript that merely restates or deepens the same package does not discharge the live observer-side bottleneck. The next honest primary task on that fixed bottleneck is therefore to prove why the package cannot be upgraded further in the required sense, which is precisely the role of a genuine necessity, maximality, or no-go theorem.
\end{proof}

\begin{proposition}[Primary-status rule for deeper-origin continuation after the second test]
\label{prop:depth}
After Theorem \ref{thm:test}, a deeper-origin note is primary only if it directly supplies one of the two missing closure-upgrade ingredients: either a stronger absorbability identity that reopens the first route on a sharpened package, or an explicit same-package field-redefinition theorem. A deeper-origin note that merely preserves the same completed observer equation without one of those upgrades is supporting context rather than the immediate next move.
\end{proposition}

\begin{proof}
The current one-metric package has now been taken through both ordered upgrade tests, and neither route is presently established there. Therefore deeper-origin continuation changes the live priority only if it directly supplies one of the specific missing upgrade ingredients rather than merely reproducing the same observer equation and the same incomplete closure status.
\end{proof}

\section{Conclusion}

The present manuscript performs the second ordered exact upgrade test on the one-metric observer-equation bottleneck. That test asks whether the completed observer equation can already be rewritten, on the currently fixed package, as exact Einstein--Hilbert plus ordinary matter through one operative metric only by means of an explicit exact invertible same-package field redefinition.

The outcome is again narrower than a no-go theorem but sharper than a generic status note. A mere relabeling of the completed metric does not furnish a new route, because it collapses back to the direct absorbability problem already tested. And the current fixed package does not yet supply the stronger ingredient that the field-redefinition route actually needs: an explicit invertible rewritten metric variable together with a proved Einstein--matter reduction theorem of the required one-metric same-package type.

This is the durable project-level result recorded here. On the current fixed package, both ordered closure-upgrade tests have now been executed and neither is presently established. The strongest proved exact observer-side theorem therefore remains effective equivalence rather than strict observer-side closure. The next honest primary paper on this same observer-side bottleneck is now a genuine necessity, maximality, or no-go theorem. A deeper-origin continuation is primary only if it directly supplies a stronger absorbability identity or an explicit same-package field-redefinition theorem.

\end{document}
